\section{The pentagon in weight sixteen modulo seventeen}
\label{sec:mod17}
The preceding character theorem uses rational compact images.  We now
compute the entire finite-field pentagon quotient in weight sixteen,
and give an integral construction that realizes every functional in that
quotient.  No equality between rational and modular pentagon kernels is
assumed.

\begin{theorem}\label{thm:pentagon17}
Over $\F_{17}$ the pentagon map in weight sixteen has the following ranks:
\[
 \operatorname{rank}P_{16}=70,\qquad
 \dim\ker P_{16}=5,\qquad
 \operatorname{rank}(P_{16}|_{[W,W]_{16}})=32.
\]
Its kernel is contained in $[W,W]_{16}$.  Consequently the induced map
\[
 W_{16}/[W,W]_{16}\longrightarrow
 \mathfrak t_{4,16}/P_{16}([W,W]_{16})
\]
is injective and has rank thirty-eight.  Every functional on
$(W/[W,W])_{16}$ annihilates the linear pentagon kernel.
\end{theorem}
\begin{proof}
We first give the upper bound on the rank.  The infinitesimal braid Lie
algebra is torsion free over $\Z$: its presentation gives the semidirect
decomposition
\[
 \mathfrak t_4=\Lie_\Z(t_{14},t_{24},t_{34})\rtimes\mathfrak t_3,
 \qquad
 \mathfrak t_3=\Z(t_{12}+t_{13}+t_{23})\oplus
                   \Lie_\Z(t_{12},t_{23}).
\]
Indeed the cross-brackets specified by the infinitesimal braid relations
define derivations of the free first factor; checking the relations on
its three generators gives inverse homomorphisms to and from this
semidirect presentation.  Each homogeneous component is therefore free
of finite rank.  Section~\ref{sec:dimensions} proves that
$W_{16}(\Z_{17})$ is free and reduces onto $W_{16}(\F_{17})$.
The rank of its integral pentagon matrix cannot increase modulo seventeen.

By the Deligne--Drinfeld theorem, the rational pentagon kernel has
dimension five in weight sixteen.  The free Lie algebra on the odd
generators has three two-letter brackets, of weights $(3,13)$, $(5,11)$
and $(7,9)$, and two four-letter words, of multidegrees $(3,3,3,7)$
and $(3,3,5,5)$.  Each of those latter multidegrees has dimension one;
this follows directly by Jacobi, or by the multigraded Witt formula.
No word with six letters has weight sixteen.  Thus the rational rank is
$75-5=70$, and the modular rank is at most seventy.

Here is a complete finite specification of a lower-bound certificate.
All row reductions in this paragraph are over $\F_{17}$, using the
first possible pivot in the stated order.  In the lexicographic Lyndon
basis of $L_{16}$ form the columns of
\[
 \Pi=(1-s)(2-r-r^2),\qquad
 s(q)=q(y,x),\quad r(q)=q(y,-x-y).
\]
The image of $\Pi$ is exactly
$\{q:sq=-q,\ (1+r+r^2)q=0\}$: on the standard representation,
$2-r-r^2$ is three times the identity, while it is zero on the other
two irreducible representations.  Retain the first independent columns
of $\Pi$, obtaining $1360$ polynomials.  Apply the special equation
\eqref{eq:special} to these columns, taking length-seventeen Lyndon-word
coefficients.  Its matrix has rank $1285$.  For each free column, in
increasing order, take the reduced-row-echelon null vector with that free
coordinate one and the other free coordinates zero.  This gives the
ordered basis $q_1,\ldots,q_{75}$ of $W_{16}$ used below.  The equations
on all rows vanish on these vectors; independence and the dimension
formula prove that no solution is omitted.

For every output word in Table~\ref{tab:rows17}, take the indicated
coefficient of $P_{16}(q)(a)$ in the three-generator representation of
Section~\ref{sec:cert}.  List the left column of words first and the
right column second, and call the resulting $70\times75$ matrix $R$.
The coefficient recurrence \eqref{eq:contraction}, split into blocks of
eight operators, gives
\begin{equation}
 \det R_{\{1,\ldots,70\},\{1,\ldots,70\}}=4\quad\text{in }\F_{17}.
 \label{eq:minor17}
\end{equation}
This proves rank at least seventy and hence equality.  In particular,
these coefficient rows have exactly the same kernel as $P_{16}$.

Let $b_1,\ldots,b_{37}$ be the ordered bracket basis of
Section~\ref{sec:brackets}, obtained by omitting the twenty-eighth of the
thirty-eight columns.  The $37\times5$ matrix $X$ in
Table~\ref{tab:kernel17} defines five vectors
$v_j=\sum_i X_{ij}b_i$.  Applying the same coefficient recurrence gives
\begin{equation}
 R(v_1\ \cdots\ v_5)=0.
 \label{eq:kernel17}
\end{equation}
The minor of $X$ on rows $1,2,3,22,30$ has determinant $5$ modulo
seventeen, so the five vectors are independent.  They therefore span
$\ker R=\ker P_{16}$ and all lie in the bracket space.  Its dimension
is thirty-seven, so the restriction of $P_{16}$ to it has rank
$37-5=32$.  The quotient assertion now follows from elementary linear
algebra.  Equations \eqref{eq:minor17}--\eqref{eq:kernel17}, together
with the deterministic basis prescription, are finite certificates of
every modular assertion used in the proof.
\end{proof}

\begin{table}[p]
\centering\small
\begin{tabular}{rllrll}\toprule
row&input&output word&row&input&output word\\\midrule
1&$a$&\texttt{abcbbaaaaaacaaaca}&36&$a$&\texttt{baacabaaaacacabba}\\
2&$a$&\texttt{ccabbbcacacbbabaa}&37&$a$&\texttt{babbaaaacbccaaaaa}\\
3&$a$&\texttt{aabaccabbbcaccabb}&38&$a$&\texttt{cabaabccabbaccbaa}\\
4&$a$&\texttt{baabcccaaaabbbaca}&39&$a$&\texttt{cbaaaccbcacbbabba}\\
5&$a$&\texttt{accbacbaacaabbacb}&40&$a$&\texttt{abaaaaabcacacabaa}\\
6&$a$&\texttt{cbcabacbabcaaabaa}&41&$a$&\texttt{aacbaaaaacaabbbac}\\
7&$a$&\texttt{cabaabaaabacaaaca}&42&$a$&\texttt{acabaacacabaaaaba}\\
8&$a$&\texttt{abaabbacaaaaaacca}&43&$a$&\texttt{abaabccacababbaca}\\
9&$a$&\texttt{aaaaabbbaaccacbbc}&44&$a$&\texttt{abacaacaabaaaacba}\\
10&$a$&\texttt{abbbacbaccacabaca}&45&$a$&\texttt{cccababcabacbabba}\\
11&$a$&\texttt{aaaaaaaacaabcbcab}&46&$a$&\texttt{aabcaaacaaaabbcaa}\\
12&$a$&\texttt{aabacccacaacbbbba}&47&$a$&\texttt{bcaaabccabbbccbaa}\\
13&$a$&\texttt{caacabcaaabcbbaaa}&48&$a$&\texttt{aabacacabaaabcaba}\\
14&$a$&\texttt{bcaaccbaabccbbbaa}&49&$a$&\texttt{aaacacbbaaacbaaaa}\\
15&$a$&\texttt{acbcaacacbabcaabb}&50&$a$&\texttt{accbaccbabcabaaba}\\
16&$a$&\texttt{bcaaccbcabaaabaaa}&51&$a$&\texttt{bbccaabacabcbbcaa}\\
17&$a$&\texttt{caabbbbcabccabcaa}&52&$a$&\texttt{abaaabcbcacacaaab}\\
18&$a$&\texttt{aaaaccaaaabacbbba}&53&$a$&\texttt{abacbbaacaacaccbb}\\
19&$a$&\texttt{aaaaaacabccabacbb}&54&$a$&\texttt{acbaacbcaabcababa}\\
20&$a$&\texttt{baaaabaaaaccbbcaa}&55&$a$&\texttt{acccbababcacabbba}\\
21&$a$&\texttt{abbaabcaaacaaacab}&56&$a$&\texttt{cacbaababbaaaacaa}\\
22&$a$&\texttt{cbbaacbaacabbcaaa}&57&$a$&\texttt{abaaaabacbbaccaaa}\\
23&$a$&\texttt{acabacbaaaacbaaaa}&58&$a$&\texttt{caabbaaaacbaaabca}\\
24&$a$&\texttt{cbbaabacbccbcabaa}&59&$a$&\texttt{aaaaaaacbacbabcaa}\\
25&$a$&\texttt{abbcbcbbacacaaabc}&60&$a$&\texttt{abaaaacbbcccbbcaa}\\
26&$a$&\texttt{bbacbacbaabbaccca}&61&$a$&\texttt{acaacabbbccbaaaaa}\\
27&$a$&\texttt{acbaccabaaabcaaba}&62&$a$&\texttt{aacbbcabcaaaaaaba}\\
28&$a$&\texttt{aabaaaaacbacabaca}&63&$a$&\texttt{aaaababaaaaabccca}\\
29&$a$&\texttt{aaaacbaabaacacbbc}&64&$a$&\texttt{bcaabbaaaaacaaaca}\\
30&$a$&\texttt{abcbaccaaababacab}&65&$a$&\texttt{aabacbccabbaaccba}\\
31&$a$&\texttt{aaaabbacbaaacaaca}&66&$a$&\texttt{cccababacabcbbaba}\\
32&$a$&\texttt{acbaabbccabacaaba}&67&$a$&\texttt{aaaacbcaabaacbaba}\\
33&$a$&\texttt{ccaaabbcaaabcbaba}&68&$a$&\texttt{aacbbaabaacaaabac}\\
34&$a$&\texttt{acaababcaaaabaaac}&69&$a$&\texttt{abacaabcbabbcacbc}\\
35&$a$&\texttt{ccaaaccaaababbbaa}&70&$a$&\texttt{aabaaaaacbacbacba}\\
\bottomrule\end{tabular}
\caption{The seventy coefficient rows of the weight-sixteen pentagon certificate over $\F_{17}$.}\label{tab:rows17}
\end{table}

\begin{table}[!htbp]
\centering\small
\begin{tabular}{r|rrrrr@{\qquad}r|rrrrr}\toprule
$i$&\multicolumn{5}{c}{$(X_{i1},\ldots,X_{i5})$}&$i$&\multicolumn{5}{c}{$(X_{i1},\ldots,X_{i5})$}\\\midrule
1&7&1&3&10&15&20&8&10&14&16&14\\
2&3&10&12&14&1&21&6&8&11&16&0\\
3&15&11&6&15&6&22&11&4&5&5&8\\
4&16&16&12&5&5&23&7&1&14&14&2\\
5&5&3&4&8&2&24&5&8&10&10&16\\
6&8&11&7&16&1&25&12&9&7&7&1\\
7&9&5&0&1&12&26&9&11&1&1&5\\
8&10&11&9&1&5&27&16&12&15&15&7\\
9&7&0&6&6&15&28&15&7&13&13&14\\
10&7&3&3&7&9&29&16&12&15&15&7\\
11&0&6&11&2&5&30&11&15&7&2&11\\
12&10&12&15&14&13&31&0&0&0&0&0\\
13&9&7&6&4&0&32&12&4&3&13&12\\
14&9&7&11&9&12&33&8&14&2&3&8\\
15&8&1&14&4&7&34&15&5&8&12&15\\
16&14&3&6&10&10&35&0&0&0&0&0\\
17&9&11&11&3&0&36&4&7&1&10&4\\
18&6&4&11&5&12&37&14&16&12&1&14\\
19&2&1&2&9&1&&&&&&\\
\bottomrule\end{tabular}
\caption{Bracket coordinates of five independent pentagon-kernel vectors. All entries are residues modulo seventeen.}\label{tab:kernel17}
\end{table}


\begin{theorem}[Integral top-weight characters]\label{thm:topweight}
For every prime $p\geq5$, put $m=p-1$.  After a choice of tame splitting,
there is an injective linear map
\[
 \bigl(W_m(\F_p)/([W,W]_m(\F_p)+\ker P_m)\bigr)^*
 \longrightarrow
 \Hom_{\cts}\bigl(\GTz/\overline{\langle\!\langle\GT\rangle\!\rangle},
                  \F_p\bigr).
\]
The character associated with $\ell$ restricts to $\ell$ on the leading
weight-$m$ subgroup.  The characters are independent already on the
global multiplier-one subgroup.  In particular the integral certificate
in weight twelve gives a nonzero character of order thirteen without
using the Deligne--Drinfeld theorem; Theorem~\ref{thm:pentagon17} realizes
all thirty-eight weight-sixteen functionals as characters of order
seventeen.
\end{theorem}
\begin{proof}
Use $G,U,P_0,P_1$ as in the proof of
Theorem~\ref{thm:even-characters}, with $m=p-1$.
The group $U$ is a finitely generated torsion-free pro-$p$ group of
nilpotency class at most $\lfloor m/3\rfloor<p$.
The Lazard correspondence therefore gives a finite free
$\Z_p$-Lie algebra $\mathcal L=\log U$, compatible with closed subgroups,
normal subgroups and automorphisms
\cite[Theorem 1(c)--(f)]{LazardTransport}.
Its filtration has associated graded $W_n(\Z_p)$ for $3\leq n\leq m$.
This can also be read directly from the finite logarithm series: its
denominators are prime to $p$, and its first term is unchanged on each
leading filtered quotient.  In particular
$\log P_1$ is a closed $\Z_p$-Lie subalgebra of $\mathcal L$.

Let $\omega\in\Z_p^\times$ generate the group of $(p-1)$st roots of
unity.  The preimage of this finite group in $P_0$ is an extension by
the pro-$p$ group $P_1$.  Choose a lift $g$ of $\omega$.  The sequence
$g^{p^j}$ converges to an element $t\in P_0$ of order $p-1$ and
multiplier $\omega$: in each finite quotient, the order of $g$ divides
$(p-1)p^a$ for some $a$, and these powers are eventually constant, being zero
modulo $p^a$ and one modulo $p-1$.  Its integral eigenspace projectors are
\[
 e_n=\frac1{p-1}\sum_{j=0}^{p-2}\omega^{-nj}\operatorname{Ad}(t)^j,
        \qquad 3\leq n\leq p-1.
\]
These weights have distinct residues modulo $p-1$, and the filtration
shows that
\[
 \mathcal L=\bigoplus_{n=3}^{m}\mathcal L^{(n)},\qquad
 \mathcal L^{(n)}=e_n\mathcal L\simeq W_n(\Z_p).
\]
Brackets respect ordinary weight: if $i+j\leq m$, their leading term
has eigenvalue $\omega^{i+j}$ and there is no other eigenspace with that
residue; if $i+j>m$, the filtration makes the bracket zero.
Thus the top component of $[\mathcal L,\mathcal L]$ is exactly the
integral Ihara bracket span.  In particular its reduction lies in
$[W,W]_m(\F_p)$, with no saturation assertion required.

For a functional $\ell$ in the theorem define
\[
 \theta_\ell(u)=\ell\bigl(e_m\log u\bmod p\bigr).
\]
The Baker--Campbell--Hausdorff corrections are brackets with
$p$-integral coefficients.  They vanish under $\ell$, so $\theta_\ell$
is a continuous homomorphism $U\to\F_p$.
Conjugation on the integral Lie abelianization factors through
$G/U=\Z_p^\times$ and commutes there with $\operatorname{Ad}(t)$.
On its top summand it acts by $\lambda_p^m$; this is one modulo $p$.
Hence $\theta_\ell$ is invariant under all of $G$.

Since $t\in P_0$, every projector $e_n$ preserves $\log P_1$.
For $h\in P_1$, the element
$\exp(e_m\log h)$ therefore belongs to $P_1$ and has no terms below
weight $m$.  Its leading term satisfies the integral linear pentagon.
Its reduction belongs to $\ker P_m$ over $\F_p$, and $\ell$ kills it.
Thus $\theta_\ell(P_1)=0$.  As $G=UP_0$, the well-defined continuous
homomorphism $\chi_\ell(up)=\theta_\ell(u)$ extends it to $G$ and
vanishes on $P_0$, exactly as in the rational proof.

The construction is linear in $\ell$.  On $U_m$ it is precisely the
prescribed functional of the leading coefficient.  Realization and
full peripheral lifting give every such coefficient on global
multiplier one, proving injectivity and the final assertions.
\end{proof}

\begin{corollary}[Strictness inside a quadratic kernel]
\label{cor:quadratic-kernel}
Let $\delta:\GTz\to\F_2$ be a continuous character.  Its kernel
contains a global multiplier-one element outside $\GT$.  If $\delta$
is nonzero and vanishes on $\GT$, then
\[
 \GT\varsubsetneqq\ker\delta\varsubsetneqq\GTz.
\]
This conclusion does not use the Deligne--Drinfeld theorem.
\end{corollary}
\begin{proof}
The weight-twelve certificate and Theorem~\ref{thm:topweight} give a
character $\chi:\GTz\to\F_{13}$ which vanishes on $\GT$ and is onto
on global multiplier one.  Choose such an element $g$ with $\chi(g)=1$.
Then $\delta(g^2)=0$ and $\chi(g^2)=2\ne0$, so $g^2$ is the required
element.  The two strict inclusions follow from the additional hypotheses
on $\delta$.
\end{proof}
